\documentclass[tikz,border=8pt]{standalone}
\usepackage{fontspec}
\usepackage{amsmath}
\usetikzlibrary{arrows.meta,decorations.markings,calc}

\setmainfont{Arial}
\setsansfont{Arial}

\definecolor{pagebg}{HTML}{F8FAFC}
\definecolor{border}{HTML}{CBD5E1}
\definecolor{textmain}{HTML}{0F172A}
\definecolor{textmuted}{HTML}{475569}
\definecolor{field}{HTML}{2563EB}
\definecolor{velocity}{HTML}{DC2626}
\definecolor{force}{HTML}{16A34A}
\definecolor{chargefill}{HTML}{FEF3C7}
\definecolor{chargeborder}{HTML}{D97706}
\definecolor{panel}{HTML}{FFFFFF}

\begin{document}
\begin{tikzpicture}[
  x=1mm,y=1mm,line cap=round,line join=round,font=\sffamily,
  vector/.style={-Latex,line width=1.2mm},
  labelbox/.style={fill=panel,draw=border,line width=0.25mm,rounded corners=1.5mm,inner sep=2.2mm}
]
  \fill[pagebg] (0,0) rectangle (230,146);
  \draw[fill=panel,draw=border,line width=0.45mm,rounded corners=4mm]
    (5,5) rectangle (225,141);

  \node[font=\sffamily\bfseries\fontsize{16}{20}\selectfont,text=textmain]
    at (115,128) {Сила Лоренца};

  \foreach \x in {32,57,82,107,132,157,182,207} {
    \foreach \y in {40,65,90} {
      \draw[field,line width=0.7mm] (\x-4,\y-4) -- (\x+4,\y+4);
      \draw[field,line width=0.7mm] (\x-4,\y+4) -- (\x+4,\y-4);
    }
  }
  \node[labelbox,font=\sffamily\fontsize{10}{12}\selectfont,text=field,align=center]
    at (48,110) {$\vec B$ от нас\\в плоскость рисунка};

  \draw[fill=chargefill,draw=chargeborder,line width=0.85mm] (94,70) circle (7.2);
  \node[font=\sffamily\bfseries\fontsize{21}{21}\selectfont,text=chargeborder]
    at (94,70) {$\boldsymbol{+}$};

  \draw[velocity,vector] (94,70) -- (151,70);
  \node[labelbox,font=\sffamily\bfseries\fontsize{12}{15}\selectfont,text=velocity]
    at (151,82) {$\vec v$};

  \draw[force,vector] (94,70) -- (94,109);
  \node[labelbox,font=\sffamily\bfseries\fontsize{12}{15}\selectfont,text=force]
    at (120,105) {$\vec F_{\text{Лор}}$};

  \draw[black,line width=0.45mm] (103,70) -- (103,81) -- (94,81);

  \node[font=\sffamily\fontsize{10}{12}\selectfont,text=textmuted,align=center]
    at (115,20) {Для отрицательного заряда направление силы меняется на противоположное.};
\end{tikzpicture}
\end{document}
